% ════════════════════════════════════════════
% 项目概述（01.引言.tex）
% ════════════════════════════════════════════
\section{项目概述}
\subsection{项目背景}

社交媒体品牌主页的帖子曝光量决定后续互动、点击与转化机会，是营销漏斗最上端的指标。品牌方需要判断什么样的内容形式、内容主题与发布时机能获得更多自然曝光，以及付费推广究竟带来多少增量。现有社交媒体实证研究多关注互动率或转发量，对曝光量的建模常停留在相关系数与分组比较，缺少对共线结构的显式处理与对预测不确定性的量化，这构成本项目的核心难点。

帖子曝光数据通常呈现三个不利于直接建模的特征。其一是账号状态与时间高度绑定：主页粉丝规模随月份单调累积，二者在样本中几乎完全共线，同时进入回归会使粉丝规模的系数失去可识别的经济含义\cite{belsley1980}。其二是曝光量为右偏计数，少数病毒式传播帖子的数值可达中位数的上百倍，直接做最小二乘会违反同方差与正态性假设，需要先做对数变换并采用稳健标准误\cite{wooldridge2019}。其三是类别分布极不均衡，视频帖与链接帖数量很少，估计精度差异很大。本项目采用覆盖数据理解、清洗整理、回归诊断、效应测算与稳健性复核的完整流程，使结论既能通过统计检验，又能转化为可执行的发帖建议。

\subsection{数据概况}

本项目的数据来自某化妆品品牌主页的帖子记录，共 500 行、8 个字段。每行对应一条帖子，字段依次为粉丝规模 X1、帖子类型 X2、内容主题 X3、发帖月份 X4、发帖星期 X5、发帖时段 X6、付费标记 X7 与曝光总次数 Y1。数据没有帖子编号与发布日期，同一主页只有一次观测，因此属于帖子级截面数据，而不是逐日汇总的时序数据。粉丝规模介于 81,370 至 139,441 人，仅有 90 个不同取值；曝光量从 570 次到 1,110,282 次，跨度超过三个数量级；全部字段中只有付费标记缺失 1 条。上述特征决定了后续必须处理重尾、类别不平衡与共线三项问题。

全部计算在 Python 3 中完成，原始文件以只读方式载入，清洗、编码与建模均基于派生的分析数据，涉及重抽样的环节固定随机种子 20261122，处理过程可由求解目录中的主程序复现。
